% ================================================================ 3.3a  Vocabulary
\begin{frame}{First, the Vocabulary: an LLM Is a Policy}
Everything that follows is just the familiar RL language, applied to text.
\vfill
\begin{itemize}\itemsep0.4em\small
    \item \textbf{Prompt $x$}, the input (a question or instruction).
          \hfill \emph{the \textbf{state}}
    \item \textbf{Response $y$}, the model's full answer.
          \hfill \emph{the \textbf{action}}
    \item \textbf{Policy $\pi_\theta(y \mid x)$}, the language model itself:
          given a prompt, it assigns a probability to possible responses and
          samples one. The weights $\theta$ are what we train.
    \item It writes the answer one \emph{token} (word-piece) at a time; at each step
          a final layer, the \textbf{LM head}, outputs a probability over the
          next token.
\end{itemize}
\vfill
\begin{center}
\empr{So ``train the policy'' simply means ``adjust the language model's weights
$\theta$.''}
\end{center}
\end{frame}

% ================================================================ 3.3  Roadmap (plain English)
\begin{frame}{The RLHF Pipeline: The Map}
Three stages. We build each one up over the next slides; for now, just the plain
English.
\vfill
\centering
\resizebox{0.98\linewidth}{!}{%
\begin{tikzpicture}[font=\footnotesize,
    box/.style={draw, rounded corners, minimum width=3.6cm, minimum height=1.7cm, align=center, fill=sysblue},
    ar/.style={-{Stealth[length=2.8mm]}, very thick}]
  \node[box] (s1) at (0, 0)  {\textbf{1.\ Imitate}\\[0.25em]\scriptsize teach it to\\\scriptsize follow instructions};
  \node[box] (s2) at (5.2, 0){\textbf{2.\ Score}\\[0.25em]\scriptsize learn what\\\scriptsize people prefer};
  \node[box] (s3) at (10.4, 0){\textbf{3.\ Optimize}\\[0.25em]\scriptsize push toward\\\scriptsize preferred answers};
  \draw[ar] (s1) -- (s2);
  \draw[ar] (s2) -- (s3);
\end{tikzpicture}%
}
\vfill
\begin{itemize}\itemsep0.25em\small
    \item \textbf{Start:} a raw language model that can produce text but doesn't
          reliably do what you ask.
    \item \textbf{End:} a model that follows instructions \emph{and} leans toward
          the answers people actually want.
\end{itemize}
\srcnote{\scriptsize Standard pipeline: Ouyang et al.\ (InstructGPT), 2022; Stiennon et al.\ 2020; Bai et al.\ 2022.}
\end{frame}

% ================================================================ 3.3 Stage 1: SFT
\begin{frame}{Stage 1 (Imitate): Supervised Fine-Tuning (SFT)}
\begin{itemize}\itemsep0.5em
    \item \textbf{Problem:} a raw pretrained model can \emph{continue} text, but
          doesn't reliably answer questions or follow instructions.
    \item \textbf{Fix: show it examples.} Collect prompts paired with
          \emph{human-written ideal answers}, and train the model to imitate them.
          This is ordinary supervised learning (next-token prediction), no RL
          yet.
    \item The result is a decent instruction-follower. We \textbf{freeze a copy of
          it} and call it the \empr{reference model $\pi_{\text{ref}}$}; it will
          be our anchor in Stages 2 and 3.
\end{itemize}
\vspace{0.5em}
\begin{center}\footnotesize
SFT can only \emph{copy} the demonstrations it was given. To go \emph{beyond} them,
we need to capture what people \emph{prefer}. That is Stage 2.
\end{center}
\end{frame}

% ---------------------------------------------------------------- 3.3 Stage 2a: the RM in words
\begin{frame}{Stage 2 (Score): A Model That Rates Answers}
\begin{itemize}\itemsep0.5em
    \item We want a \textbf{reward model} $r_\phi$: a network (its own weights
          $\phi$) that reads a (prompt, response) and outputs \emph{a single
          number} (\emph{how good is this answer?}). Higher means better.
    \item \textbf{Where does its training data come from?} Comparisons. For one
          prompt $x$, show a human two candidate responses; they pick the better
          one. We label the pair:
    \begin{itemize}\footnotesize\itemsep0.15em
        \item $y^+$, the response the human \textbf{preferred},
        \item $y^-$, the response they \textbf{rejected}.
    \end{itemize}
    \item \textbf{What we want from $r_\phi$:} give $y^+$ a higher score than $y^-$,
          on every pair. Next slide: how to turn ``preferred'' into a trainable
          loss.
\end{itemize}
\end{frame}

% ---------------------------------------------------------------- 3.3 Stage 2b: Bradley-Terry loss
\begin{frame}{Turning Comparisons into a Loss: Bradley--Terry}
How likely is a person to prefer $y^+$ over $y^-$? It should grow with the
\textbf{gap in their scores}, squeezed into a probability by the sigmoid $\sigma$
(``$\succ$'' reads \emph{``is preferred over''}):
\[
    \underbrace{P\big(y^+ \succ y^- \mid x\big)}_{\text{chance the human picks } y^+}
    \;=\; \sigma\big(\, r(y^+) - r(y^-) \, \big),
    \qquad \sigma(z) = \frac{1}{1 + e^{-z}}
\]
Train $r_\phi$ to make these probabilities match the human labels (maximum
likelihood). The negative log-likelihood is the \textbf{reward-model loss}:
\[
    L_{\mathrm{RM}} \;=\; -\log\sigma\big(\, r_\phi(y^+) - r_\phi(y^-) \, \big)
\]
\begin{itemize}\itemsep0.2em\small
    \item In plain words: \empr{push the preferred answer's score above the rejected
          one's.}
    \item Remember this $-\log\sigma(\Delta)$ shape: the \textbf{exact same form}
          comes back when we reach DPO.
    \item Only the \emph{gap} between scores matters; the absolute scale is free.
\end{itemize}
\end{frame}

